\section{条件生成把控制信号写进得分}
\label{sec:14-Deep-Learning/08-Diffusion-Models/04}

你手上有一个训好的无条件模型，能生成各种各样的鸟。产品要的是``只出红色的鸟''。重新训一个只喂红鸟数据的模型是最直接的想法，但你手上的红鸟图片只有几千张，训出来的东西质量掉一大截。退而求其次，照常采样然后筛掉不合要求的，一百张里能挑出三张，成本涨三十倍，而且用户每换一个要求就得重来一遍。

这两条路都在把条件当成\emph{外部}的东西——要么在训练前筛数据，要么在生成后筛结果。真正管用的做法是把条件写进采样过程本身，让每一步都朝着满足条件的方向偏一点。而写进去的位置，恰好就是前两节反复出现的那个量。

\subsection{支点：条件得分等于无条件得分加一项修正}

从 Bayes 公式开始，取对数：

\[
\log p(x\given y)=\log p(x)+\log p(y\given x)-\log p(y).
\]

对 \(x\) 求梯度。\(\log p(y)\) 不含 \(x\)，导数为零，直接消失：

\[
\nabla_x\log p(x\given y)=\nabla_x\log p(x)+\nabla_x\log p(y\given x).
\]

推导只有两行，但值得盯一会儿。左边是条件分布的得分——如果有了它，把上一节的采样算法里的得分换掉，跑出来的就是符合条件 \(y\) 的样本。右边把它拆成了两块：一块是原来那个无条件模型已经会的东西，另一块是一个\emph{判别式}的量，即``给定这张图，它有多像红色的鸟''对图像的梯度。条件生成于是变成了在原有得分场上叠加一个修正向量场。

那个消失的 \(\log p(y)\) 值得单独记一笔。在密度这一层，条件分布 \(p(x\given y)=p(x)p(y\given x)/p(y)\) 里的归一化常数需要一个高维积分才能算，这是整本书里反复出现的那个障碍。换到得分这一层，它是常数，一求导就没了。用得分而不用密度来表示分布，买到的正是这一点：所有不依赖 \(x\) 的归一化因子自动退场。\hyperref[sec:06-Random-Processes/06-Diffusion-and-SDE/05]{``得分匹配：怎样从样本学习得分''}里用同一个理由绕开了配分函数，这里用它绕开的是条件化的归一化常数。

\subsection{用分类器提供那一项，和它的两个麻烦}

按这个式子直接施工：训一个分类器 \(p_\phi(y\given x_t,t)\)，采样时把它对输入的梯度加到得分上。上一节里得分和噪声预测差一个已知缩放 \(\nabla_x\log p_t(x)=-\varepsilon_\theta(x,t)/\sqrt{1-\bar\alpha_t}\)，换算过去就是把网络输出改成

\[
\hat\varepsilon(x_t,t)=\varepsilon_\theta(x_t,t)-s\,\sqrt{1-\bar\alpha_t}\;\nabla_{x_t}\log p_\phi(y\given x_t,t),
\]

其中 \(s\) 是一个人为放大倍数，\(s=1\) 对应上面那个恒等式。这就是 classifier guidance。它有个真正的好处：无条件模型不用重训，任何新的条件只要配一个新分类器。

麻烦有两个，都不小。分类器必须在\emph{加噪}样本上训练——采样中途的 \(x_t\) 是一张被噪声盖住的图，你从别处拿来的、在干净图上训的现成分类器喂进去输出的是垃圾，而 \(t\) 大的时候图上几乎什么都没有，这个分类任务本身就很难。第二个麻烦更隐蔽：采样过程在沿着分类器的梯度上升，这正是生成对抗样本的做法。模型完全可以找到一些让分类器给出极高置信、但人看着根本不像红鸟的方向，然后一路走过去。你以为在优化``像红鸟''，实际在优化``骗过这个分类器''，两者只在分类器可靠的区域里一致。

\subsection{不用分类器也能拿到那一项}

把支点等式移一下项：

\[
\nabla_x\log p(y\given x)=\nabla_x\log p(x\given y)-\nabla_x\log p(x).
\]

需要的那个修正项，本身就等于条件得分与无条件得分之差。如果一个模型同时会算这两个得分，分类器就不必存在了。

做法朴素得很：训练时以某个概率（通常一两成）把条件 \(y\) 替换成一个专门的空标记 \(\varnothing\)，别的什么都不改。同一组参数于是既学会了 \(\varepsilon_\theta(x_t,y)\)，也学会了 \(\varepsilon_\theta(x_t,\varnothing)\)。采样时把两者做一次外插：

\[
\tilde\varepsilon=\varepsilon_\theta(x_t,\varnothing)+w\bigl(\varepsilon_\theta(x_t,y)-\varepsilon_\theta(x_t,\varnothing)\bigr).
\]

这就是 classifier-free guidance。\(w=0\) 时退回无条件生成；\(w=1\) 时右边化简成 \(\varepsilon_\theta(x_t,y)\)，也就是老老实实按条件得分采样；\(w>1\) 时把两者的差放大，往``更条件''的方向过冲。换来这份方向感的是每一步要跑两次网络（一次带条件、一次不带），采样成本翻倍。

\subsection{\texorpdfstring{\(w>1\) 时你采样的不是条件分布}{w>1 时你采样的不是条件分布}}

把 \(\tilde\varepsilon\) 换算回得分，用一次支点等式：

\[
\begin{aligned}
\text{等效得分}
&=\nabla_x\log p(x)+w\bigl(\nabla_x\log p(x\given y)-\nabla_x\log p(x)\bigr)\\
&=\nabla_x\log p(x)+w\,\nabla_x\log p(y\given x)\\
&=\nabla_x\log\Bigl(p(x)\,p(y\given x)^{w}\Bigr).
\end{aligned}
\]

所以沿这个场采样，采到的分布正比于 \(p(x)\,p(y\given x)^{w}\)。\(w=1\) 时它就是 \(p(x\given y)\)；\(w>1\) 时它\emph{不是}任何条件分布，而是把似然项抬到 \(w\) 次幂之后重新归一化的结果。

指数放大对密度的作用是锐化。\(p(y\given x)\) 高的地方被相对抬得更高，低的地方被压得更低，倍率随 \(w\) 指数增长。如果条件分布本来有几个模式，其中一个的条件似然稍高一点，取幂之后这点优势被放大，其余模式的质量占比迅速塌下去。

\begin{center}
\begin{tikzpicture}[font=\small,>=stealth]
% ================= 左：两条反向曲线 =================
\begin{scope}
  \fill[orange!12] (1.08,0) rectangle (4.32,2.75);
  \draw[orange!55,dashed] (1.08,0) -- (1.08,2.75);
  \draw[orange!55,dashed] (4.32,0) -- (4.32,2.75);
  \node[font=\footnotesize,orange!70!black] at (2.70,2.98) {常用工作区间};
  \draw[->,black!55] (-0.15,0) -- (5.8,0) node[right,font=\footnotesize] {\(w\)};
  \draw[->,black!55] (-0.15,0) -- (-0.15,3.2);
  \foreach \p/\lab in {0/0, 1.08/2, 2.16/4, 3.24/6, 4.32/8, 5.4/10}{
    \draw[black!45] (\p,0) -- (\p,-0.12);
    \node[below,font=\footnotesize] at (\p,-0.12) {\lab};
  }
  \draw[blue!70,thick,smooth] plot coordinates
    {(0,0.30)(0.54,1.05)(1.08,1.62)(1.62,1.95)(2.16,2.14)(3.24,2.32)(4.32,2.40)(5.4,2.44)};
  \draw[red!70,thick,densely dashed,smooth] plot coordinates
    {(0,2.45)(0.54,2.30)(1.08,2.05)(1.62,1.75)(2.16,1.45)(3.24,0.95)(4.32,0.62)(5.4,0.42)};
  \node[right,font=\footnotesize,blue!70] at (2.55,1.80) {与提示的贴合度};
  \node[right,font=\footnotesize,red!70] at (3.30,0.38) {样本多样性};
  \node[below,font=\footnotesize] at (2.7,-0.62) {旋钮：\(w\) 调大一格，两边同时动};
\end{scope}
% ================= 右：密度被取幂锐化 =================
\begin{scope}[shift={(7.1,0)}]
  \draw[->,black!55] (-0.15,0) -- (3.75,0) node[right,font=\footnotesize] {\(x\)};
  \draw[->,black!55] (-0.15,0) -- (-0.15,3.2);
  \draw[black!45,thick,smooth] plot coordinates
    {(0,0.05)(0.3,0.20)(0.6,0.48)(0.9,0.70)(1.1,0.75)(1.4,0.62)(1.7,0.42)
     (1.9,0.38)(2.2,0.52)(2.5,0.60)(2.8,0.45)(3.1,0.22)(3.4,0.07)};
  \draw[violet!65,thick,smooth] plot coordinates
    {(0,0.01)(0.3,0.05)(0.6,0.35)(0.9,1.10)(1.1,1.35)(1.4,0.85)(1.7,0.28)
     (1.9,0.20)(2.2,0.38)(2.5,0.62)(2.8,0.32)(3.1,0.07)(3.4,0.01)};
  \draw[red!70,thick,smooth] plot coordinates
    {(0,0.00)(0.3,0.00)(0.6,0.10)(0.9,1.55)(1.1,2.35)(1.4,1.20)(1.7,0.10)
     (1.9,0.04)(2.2,0.10)(2.5,0.30)(2.8,0.08)(3.1,0.01)(3.4,0.00)};
  \node[right,font=\footnotesize,black!55] at (2.55,0.72) {\(w=1\)};
  \node[right,font=\footnotesize,violet!65] at (1.28,1.52) {\(w=3\)};
  \node[right,font=\footnotesize,red!70] at (1.18,2.42) {\(w=8\)};
  \node[below,font=\footnotesize] at (1.7,-0.62) {\(p(x)\,p(y\given x)^{w}\) 随 \(w\) 变尖};
\end{scope}
\end{tikzpicture}

\emph{左边是同一个旋钮的两侧读数：往右拧，样本更听话，也更雷同。右边给出原因——取 \(w\) 次幂让密度整体收窄，次要模式的质量被抽干，最后几乎所有样本都落在同一个峰上。要记住的是这两幅图讲的是同一件事，不是两个独立的现象。}

\end{center}

图右边那个次要的峰被抽干，落到实际使用上就是：同一句提示换不同随机种子，\(w\) 小的时候能出各种构图，\(w\) 大的时候出来的图彼此高度相似。再大一些还会出现过饱和的色块和奇怪的高对比边缘——那是外插把 \(\tilde\varepsilon\) 的幅度推到了训练时从未见过的范围，网络在分布外的地方给出的输出不再可靠。常用区间落在图上那一段，具体到哪个数取决于模型、数据和提示的类型，是调出来的，不是算出来的。

这条``贴合度换多样性''的曲线不是扩散模型独有。\hyperref[sec:14-Deep-Learning/07-Generative-Models-Unified/01]{``四条路线在最小化哪个散度''}把生成模型按它们优化的散度排过队，那里的分界是：有的目标惩罚``该覆盖的没覆盖''，逼模型把所有模式都照顾到，于是往低密度区涂抹；有的目标惩罚``不该有的地方有质量''，逼模型缩到高密度区里，于是丢掉模式。把 \(w\) 从 \(1\) 往上调，做的正是沿这条轴往后者移动——训练目标没变，但采样时实际对齐的分布换了一个，换到了一个更集中、更保守的版本上。同一条交易，换了个位置出现。

\subsection{一次难的匹配，拆成很多次容易的匹配}

回头看这四节做的事。前向过程用一条手工写死的路径把数据推成噪声，代价是零，因为热传导的抹平行为不需要学。训练把变分下界拆成一串高斯之间的 KL，每一项都塌成平方距离，最后剩下一个跨所有噪声水平的去噪回归——生成建模里最难的那类问题被写成了监督学习里最简单的那个目标。采样把学到的得分场积分回去，问题于是变成一道数值分析题，步数就是精度的价钱。条件生成在得分上加一项修正，因为得分表示让归一化常数自动消失。

贯穿始终的是同一招：不去一次性地把一个复杂分布对齐到数据，而是把这段距离切成很多小段，每一段短到可以用高斯近似、可以闭式求解、可以用一个网络一次前向搞定。切得越细，每一步越容易，这是模型好训的全部原因。

而账要算在另一边。切得越细，走完全程要走的步数越多，这是模型采样慢的全部原因。扩散模型的代价结构因此非常清晰：训练便宜——就是回归，损失函数一行，没有对抗博弈，没有配分函数，收敛稳定得让人放心；采样昂贵——要跑几十上百次网络前向，而且这个成本每生成一张图都要付一次。GAN 恰好是镜像的：训练是两个网络的博弈，不收敛、模式坍缩、超参数敏感，能训出来的人不多；一旦训出来，生成一个样本就是一次前向，几毫秒。

把难度从训练搬到采样，还是从采样搬到训练，是这两类模型之间真正的分歧，也是过去几年里扩散模型胜出的原因——训练的痛苦要研究者承担，采样的成本可以靠工程和硬件摊掉，而蒸馏、高阶求解器、隐空间扩散都是在往回搬这部分成本。没有哪一种做法能同时便宜训练和便宜采样：那个复杂分布的难度守恒，你只能选择在哪里付账。
